%_____________________________________________________________________________________________ 
% Chapter 3: Related Work
%_____________________________________________________________________________________________ 
\chapter{Related Work}
\label{sec:related_work}

It helps to understand what each related system does well and where it
stops short, since our design borrows selectively from several of them.

\section{BFT Consensus Protocols}

Practical Byzantine Fault Tolerance (PBFT)~\cite{pbft1999} proved that
practical BFT was possible, but its $O(n^2)$ all-to-all communication
makes it painful beyond a couple of dozen nodes.
\hotstuff{}~\cite{hotstuff2019} brought that down to $O(n)$ by
funnelling everything through a leader and using QCs.  More recent
variants such as HotStuff-1~\cite{hotstuff1_2024} try to cut one
communication phase, and Sync HotStuff~\cite{synchotstuff2020} gets
better throughput when the network is synchronous, focusing on
squeezing out better latency or throughput from the core protocol.
None of them tackle the four practical gaps we are interested in.

\section{Persistence}

For persistence, the closest analogue is Raft~\cite{raft2014}.  Its
WAL-plus-snapshot approach has been tested in production systems such
as etcd, and the recovery semantics are well understood.  The catch is
that Raft assumes only crash faults; transplanting its storage patterns
into a Byzantine setting requires some care.  For instance, a
recovering replica cannot just trust its own log---it needs checksum
verification and must re-validate against the quorum.

\section{Dynamic Membership}

DynaNet~\cite{dynanet2025} is one of the few BFT systems that
explicitly supports dynamic validator sets.  Our membership module
takes a simpler route: we treat \texttt{JOIN} and \texttt{REMOVE} as
regular consensus commands so that the existing 3-chain commit rule
provides the atomicity guarantee for free.

\section{Pluggable State Machines}

Tendermint~\cite{tendermint} deserves credit for popularising the idea
of separating consensus from application logic through its ABCI.
However, its main focus is on decentralised applications (dApps).  Our
\texttt{App} trait and gRPC interface borrow this philosophy directly
and apply it to applications in general.  We have targeted a
Rust-native trait as the primary interface, with gRPC as an optional
bridge for applications written in other languages.

\section{Liveness under Low Load}

As for liveness, the idea of injecting no-op commands when the system
is idle goes back to early Paxos literature.  We are not claiming
novelty there, just noting that nobody seems to have wired it into a
HotStuff pacemaker properly.
